\chapter{Methodology}
\label{ch:methodology}

\section{Introduction}
This chapter details the experimental methodology used to evaluate specific XAI techniques on tabular data. We focus on the Adult Income dataset and employ two popular explainability methods: SHAP and LIME.

\section{Experimental Design}

\subsection{Dataset}
We utilize the UCI Adult Income dataset~\cite{kohavi1996adult}.
\begin{itemize}
    \item \textbf{Total Samples}: 48790
    \item \textbf{Training Samples}: 39032
    \item \textbf{Test Samples}: 9758
    \item \textbf{Features}: 108
    \item \textbf{Positive Class Ratio}: 0.24
\end{itemize}

\subsection{Preprocessing}
[Describe imputation, encoding, and scaling steps here]

\section{Model Architectures}
We evaluate two primary black-box models: Random Forest~\cite{breiman2001random} and XGBoost~\cite{chen2016xgboost}. The specific hyperparameters used are detailed in Table~\ref{tab:hyperparameters}.

\begin{table}[htbp]
\centering
\caption{Model Hyperparameters}
\label{tab:hyperparameters}
\begin{tabular}{lll}
\toprule
Model & Parameter & Value \\
\midrule
xgboost & n\_estimators & 100 \\
xgboost & max\_depth & 6 \\
xgboost & learning\_rate & 0.1 \\
xgboost & random\_state & 42 \\
xgboost & n\_estimators & 100 \\
xgboost & max\_depth & 6 \\
xgboost & learning\_rate & 0.1 \\
xgboost & random\_state & 42 \\
rf & n\_estimators & 200 \\
rf & max\_depth & None \\
rf & random\_state & 42 \\
rf & criterion & gini \\
rf & min\_samples\_leaf & 1 \\
rf & n\_estimators & 200 \\
rf & max\_depth & None \\
rf & random\_state & 42 \\
rf & criterion & gini \\
rf & min\_samples\_leaf & 1 \\
\bottomrule
\end{tabular}
\end{table}

\section{Explainability Techniques}
We evaluate the following local explanation methods:
\begin{itemize}
    \item \textbf{SHAP}~\cite{lundberg2017shap}: Using TreeExplainer.
    \item \textbf{LIME}~\cite{ribeiro2016lime}: Using LimeTabularExplainer.
\end{itemize}

\section{Evaluation Metrics Framework}
We employ a diverse set of quantitative metrics to assess explanation quality. A summary is provided in Table~\ref{tab:metrics}.

\begin{table}[htbp]
\centering
\caption{Evaluation Metrics Summary}
\label{tab:metrics}
\begin{tabular}{lp{10cm}}
\toprule
Metric & Description \\
\midrule
Fidelity & Measures how well the explanation model mimics the black-box model locally. \\
Stability & Measures how much the explanation changes when the input/sampling is perturbed. \\
Sparsity & Measures the sparsity/complexity of an explanation. \\
\bottomrule
\end{tabular}
\end{table}

\subsection{Statistical Analysis}
To validate our findings, we apply the Friedman test and Wilcoxon signed-rank test~\cite{demsar2006statistical}.

\section{Reproducibility}
All experiments were conducted using the following environment:
\begin{itemize}
    \item Python: 3.13.2
    \item Scikit-learn: 1.8.0
    \item SHAP: 0.50.0
    \item LIME: 0.2.0.1
    \item XGBoost: 3.1.2
\end{itemize}

